% !TEX program = xelatex
\documentclass[841mm,1189mm]{tikzposter}
\usepackage{fontspec}
\usepackage{xeCJK}
\usepackage{polyglossia}
\usepackage{qrcode}
\usepackage{graphicx}
\setdefaultlanguage{vietnamese}

% Fonts
\setmainfont{Linux Biolinum O}[BoldFont=Linux Biolinum O Bold]
\setsansfont{Linux Biolinum O}[BoldFont=Linux Biolinum O Bold]

% Colors
\definecolor{iEEAIBlue}{RGB}{0,51,102}
\definecolor{iEEAIGold}{RGB}{212,175,55}

% Poster style
\usetheme{Basic}
\usecolorstyle{Default}
\title{}
\author{}
\date{}

\begin{document}

% ===== HEADER: dùng ảnh nền chính thức từ BTC =====
\node[anchor=north] at (0, 58cm) {
  \includegraphics[width=0.95\paperwidth]{poster_bg.png}
};

\vspace{12cm}

% ===== TIÊU ĐỀ SẢN PHẨM =====
{\fontsize{44}{52}\selectfont\bfseries\color{iEEAIBlue} XÂY DỰNG PHÒNG THÍ NGHIỆM AI OFFLINE (AINATOMY)}

\vspace{0.5cm}

{\fontsize{28}{36}\selectfont\color{iEEAIBlue} Phát triển năng lực kiểm chứng AI cho học sinh THPT theo CT GDPT 2018}

\vspace{0.8cm}

% ===== THÔNG TIN TÁC GIẢ =====
{\Large\bfseries Giáo viên thực hiện:} {\Large Mai Xuân Văn}

{\Large\bfseries Đơn vị công tác:} {\Large Trường THPT Thuận Hoá, Đại học Sư Phạm, Đại học Huế}

\vspace{1cm}

% ===== 5 PHẦN NỘI DUNG =====
\begin{columns}
\column{0.5}

\block{1. Bối cảnh và vấn đề cần giải quyết}{
  \Large
  \textbf{Chính sách:} QĐ 2422/QĐ-BGDĐT yêu cầu 12 tiết/năm cho giáo dục AI trong THPT.
  
  \textbf{Hạ tầng:} Phòng máy yếu, internet không ổn định ở vùng khó khăn.
  
  \textbf{Năng lực:} Học sinh thiếu kỹ năng kiểm chứng đầu ra AI.
  
  \textbf{Vấn đề:} Làm thế nào để học sinh THPT hiểu, kiểm tra và bắt lỗi AI khi hạ tầng công nghệ hạn chế?
}

\block{2. Giải pháp và công cụ AI được sử dụng}{
  \Large
  \textbf{AInatomy} -- Phòng thí nghiệm AI offline giúp học sinh quan sát quy trình AI từ đầu đến cuối.
  
  \textbf{Đặc điểm:}
  \begin{itemize}
    \item Offline hoàn toàn, không cần internet/tài khoản
    \item Tương thích máy cấu hình thấp (RAM 2GB)
    \item 28 bài học theo 22 YCCĐ Tin học 10
    \item Mô phỏng 3D bằng Three.js
  \end{itemize}
  
  \textbf{Công cụ AI:} Three.js (mô phỏng 3D), JavaScript (logic), HTML5/CSS3 (giao diện).
}

\column{0.5}

\block{3. Quá trình triển khai trong thực tế}{
  \Large
  \textbf{Giai đoạn 1: Thiết kế} (3 tháng)
  
  Phân tích CT GDPT 2018, thiết kế kiến trúc, xây dựng prototype.
  
  \textbf{Giai đoạn 2: Phát triển} (6 tháng)
  
  28 bài học theo 22 YCCĐ, mô phỏng AI 3D, hệ thống chấm điểm tự động.
  
  \textbf{Giai đoạn 3: Đánh giá} (3 tháng)
  
  Triển khai tại THPT Thuận Hoá, thu thập dữ liệu, đánh giá năng lực.
}

\block{4. Kết quả và tác động đạt được}{
  \Huge
  \textbf{28/28} YCCĐ Tin học 10
  
  \textbf{97/97} kiểm thử tự động đạt
  
  \textbf{100\%} offline
  
  \textbf{3D} mô phỏng Three.js
  
  \vspace{0.3cm}
  
  \Large
  \textbf{Tác động:} Học sinh phát triển năng lực kiểm chứng AI, hiểu quy trình AI, tự tin sử dụng AI có trách nhiệm. Giáo viên có công cụ dạy AI hiệu quả.
}

\end{columns}

\vspace{0.5cm}

\block{5. Tính bền vững và khả năng mở rộng}{
  \Large
  \textbf{Bền vững:} Offline, mã nguồn mở, tương thích máy cấu hình thấp.
  
  \textbf{Nhân rộng:} Mọi trường THPT triển khai được, không cần đào tạo phức tạp.
  
  \textbf{Phát triển:} Thêm bài học nâng cao, tích hợp AI thực tế, nghiên cứu tác động dài hạn.
}

\vspace{0.5cm}

% ===== QR CODE =====
{\Huge\bfseries\color{iEEAIBlue} Trải nghiệm AInatomy}

\vspace{0.3cm}

\qrcode[height=5cm]{https://mxuanvan02.github.io/ainatomy/}

\vspace{0.2cm}

{\Large\texttt{https://mxuanvan02.github.io/ainatomy/}}

\end{document}
